% ==========================
% 卷面变量设置（只需改这里）
% ==========================
\newcommand{\paperlevel}{初中一}
\newcommand{\papersubject}{数学 试卷一}
\newcommand{\schoolname}{金宝培元独立中学}
\newcommand{\examname}{年终考}
\newcommand{\classinfo}{J1(H / I/ F)}
\newcommand{\examdate}{30-10-2026}
\newcommand{\examtime}{13:00-14:15}
\newcommand{\questiontotal}{6}
\newcommand{\paperpercentage}{60\%}
\newcommand{\setterteacher}{蔡杰庆师}
\newcommand{\reviewteacher}{禤佩娟师}
\newcommand{\secondreviewteacher}{张惠雯师}
\newcommand{\auditteacher}{雷泳健师}



\usepackage{xeCJK}
\setCJKmainfont{NSimSun}
\newCJKfontfamily\heiti{NSimSun}
\usepackage{geometry}
\usepackage{array}
\usepackage{tikz}
\usepackage{amsmath}
\usepackage{wrapfig}
\usepackage{adjustbox}
\usepackage{enumitem}
\usepackage{fancyhdr}
\usetikzlibrary{quotes,angles}
\usetikzlibrary{calc}
\usepackage{scalerel}
\def\mcirc{\mathbin{\scalerel*{\circ}{j}}}
\def\msquare{\mathord{\scalerel*{\Box}{gX}}}

% \renewcommand{\choiceshook}{%
%     \setlength{\leftmargin}{18pt}%
% }
\geometry{margin=2cm}
\AtBeginDocument{\zihao{4}}


%region ----------- 调整 exam -----------
\examsetup{
    % problem/label = \arabic*., 
    % 设置大题题号格式 
        problem = {
        top-sep = 1em,      % 题头上方的间距（默认 0.25em）
        bottom-sep = 1em,   % 题号下方的间距（默认 0pt）
    },
    % 隐藏大题后的分值显示 
    problem/show-points = false,
    % 调整各 question 之间的垂直间距
    question = {
        top-sep = 1em,
        bottom-sep = 1em,
    },
    % 调整小题列表的缩进和间距
    list/step-name = {}, 
    solution = {
    show-solution = show-move, % 答案自动移动到卷末
    label-content = {参考答案},  % 答案区的标题文字
    show-qed = true            % 在解答末尾显示结束符号 ◻ 
  },
    page = {
    size = a4paper,           % 设定纸张大小 
    foot-content = { - ; - }  % 页脚显示为 -1-, -2- 格式 
  }
}
\setlist[enumerate,1]{
    label=(\alpha*), 
    itemsep=0.5em, 
    leftmargin=2em, 
    after=\hfill % 配合您之前要求的百分比靠右样式
}
%endregion

\newcommand{\qscoreoffset}{4em}
\newlist{questionparts}{enumerate}{1}
\setlist[questionparts]{
    label=(\alph*),
    itemsep=0.5em,
    leftmargin=2em,
    rightmargin=4em,
    before={\renewcommand{\qscoreoffset}{8em}},
}
\newcommand{\qscore}[2][\%]{%
    \leavevmode\unskip\nobreak\hfill
    \makebox[0pt][l]{\makebox[\qscoreoffset][r]{(#2#1)}}%
}

\newcommand{\makeexamheader}{%
%region 三栏标题区
% ==========================
\noindent
\begin{tabular*}{\textwidth}{@{\extracolsep{\fill}} p{0.3\textwidth} p{0.35\textwidth} p{0.3\textwidth}}

\raggedright
\textbf{\paperlevel}\\[0.2em]
\textbf{}{\papersubject}\\[0.4em]
姓名：\underline{\hspace{3cm}} &

\centering
\textbf{\schoolname}\\[0.2em]
\textbf{\examname}\\[0.4em]
班级：\classinfo &

\raggedright
日期：\examdate\\[0.2em]
时间：\examtime\\[0.4em]
得分：\underline{\hspace{2.5cm}} \\

\end{tabular*}

\begin{flushleft}
作答题：\questiontotal 题全做（\paperpercentage） 
\end{flushleft}
%endregion
}